\documentclass[lang=cn,10pt]{elegantbook}

\title{给孩子的语文书}
\subtitle{LaTeX 经典之作}

\author{齐齐爸}
\institute{共同进步组织}
\date{June 13, 2026}
\version{0.1}
\bioinfo{自定义}{英文寄语}

\extrainfo{不要以为抹消过去，重新来过，即可发生什么改变。—— 比企谷八幡}

\setcounter{tocdepth}{3}

\logo{logo-blue.png}
\cover{cover.jpg}

% 本文档命令
\usepackage{array}
\newcommand{\ccr}[1]{\makecell{{\color{#1}\rule{1cm}{1cm}}}}

% 修改标题页的橙色带
% \definecolor{customcolor}{RGB}{32,178,170}
% \colorlet{coverlinecolor}{customcolor}

\begin{document}

\maketitle
\frontmatter

\tableofcontents

\mainmatter

\chapter{句式改写练习}

逐句给出 FCE 级别参考改写 + 亮点解析，帮你建立“从普通到高级”的语感。Rewrite the sentences using more advanced vocabulary and structures
\begin{itemize}
  \item More advanced vocabulary
  \item Complex sentence structures (clauses, participles, passive voice)
  \item More formal linking words
\end{itemize}


\section{6月12日}




\bigskip

\begin{enumerate}[leftmargin=2em, label=\arabic*.]
    \item I think social media is bad for teenagers. \\[2ex]
    \rule{\linewidth}{0.4pt}
    

    \item Many people use cars, so there is a lot of pollution. \\[2ex]
    \rule{\linewidth}{0.4pt}
    

    \item We should do more exercise because it is good for our health. \\[2ex]
    \rule{\linewidth}{0.4pt}
    

    \item The film was good, but it was too long. \\[2ex]
    \rule{\linewidth}{0.4pt}
    

    \item If you don’t study, you will fail the exam. \\[2ex]
    \rule{\linewidth}{0.4pt}
    

    \item In my opinion, schools should give less homework. \\[2ex]
    \rule{\linewidth}{0.4pt}
    

    \item Tourism helps the economy, but it can damage the environment. \\[2ex]
    \rule{\linewidth}{0.4pt}
    

    \item The government must do something about climate change. \\[2ex]
    \rule{\linewidth}{0.4pt}
    

    \item I like reading books more than watching TV. \\[2ex]
    \rule{\linewidth}{0.4pt}
    

    \item She didn’t get the job because she was late for the interview. \\[2ex]
    \rule{\linewidth}{0.4pt}
    
\end{enumerate}


\section{Model Answers and Key Expressions}

\begin{enumerate}[leftmargin=2em, label=\arabic*.]

    \item \textbf{Original:} I think social media is bad for teenagers. \\
    \textbf{Rewrite:} From my perspective, social media can have a negative impact on teenagers’ well-being. \\
    \textit{Notes:} “From my perspective” is more formal than “I think”; “have a negative impact on” is more academic than “is bad for”; “well-being” covers mental and emotional health.

    \vspace{2ex}
    \hrule
    \vspace{2ex}

    \item \textbf{Original:} Many people use cars, so there is a lot of pollution. \\
    \textbf{Rewrite:} Due to the widespread use of private vehicles, air pollution levels have risen significantly. \\
    \textit{Notes:} “Due to” introduces cause more formally than “so”; “widespread use of private vehicles” is more specific; “levels have risen significantly” is more precise.

    \vspace{2ex}
    \hrule
    \vspace{2ex}

    \item \textbf{Original:} We should do more exercise because it is good for our health. \\
    \textbf{Rewrite:} It is highly advisable to engage in regular physical activity, as it contributes to overall health. \\
    \textit{Notes:} “It is highly advisable” is more objective; “engage in regular physical activity” is more formal; “contributes to overall health” adds depth.

    \vspace{2ex}
    \hrule
    \vspace{2ex}

    \item \textbf{Original:} The film was good, but it was too long. \\
    \textbf{Rewrite:} Although the film was enjoyable, it was somewhat overlong. \\
    \textit{Notes:} “Although” shows complex structure; “enjoyable” is more descriptive; “somewhat overlong” is more nuanced.

    \vspace{2ex}
    \hrule
    \vspace{2ex}

    \item \textbf{Original:} If you don’t study, you will fail the exam. \\
    \textbf{Rewrite:} Unless you study consistently, you are likely to fail the exam. \\
    \textit{Notes:} “Unless” replaces “if not”; “consistently” adds detail; “are likely to” is more realistic than “will”.

    \vspace{2ex}
    \hrule
    \vspace{2ex}

    \item \textbf{Original:} In my opinion, schools should give less homework. \\
    \textbf{Rewrite:} Personally, I believe that schools ought to reduce the amount of homework assigned to students. \\
    \textit{Notes:} “Ought to” adds strength; “reduce the amount of homework assigned” is more formal and precise.

    \vspace{2ex}
    \hrule
    \vspace{2ex}

    \item \textbf{Original:} Tourism helps the economy, but it can damage the environment. \\
    \textbf{Rewrite:} While tourism boosts the local economy, it can also have an adverse effect on the environment. \\
    \textit{Notes:} “While” contrasts more formally; “boosts” is more dynamic; “adverse effect” is academic.

    \vspace{2ex}
    \hrule
    \vspace{2ex}

    \item \textbf{Original:} The government must do something about climate change. \\
    \textbf{Rewrite:} It is essential that the government take immediate action to tackle climate change. \\
    \textit{Notes:} “It is essential that…” is stronger and more formal; “tackle” is a key FCE verb.

    \vspace{2ex}
    \hrule
    \vspace{2ex}

    \item \textbf{Original:} I like reading books more than watching TV. \\
    \textbf{Rewrite:} I much prefer reading books to watching television. \\
    \textit{Notes:} “Prefer A to B” is a fixed structure; “much prefer” intensifies meaning concisely.

    \vspace{2ex}
    \hrule
    \vspace{2ex}

    \item \textbf{Original:} She didn’t get the job because she was late for the interview. \\
    \textbf{Rewrite:} Her lateness for the interview resulted in her not being offered the job. \\
    \textit{Notes:} “Resulted in” expresses cause–effect formally; gerund passive raises grammatical level.

\end{enumerate}

\section{Useful Strategies for FCE Writing}

\begin{itemize}[leftmargin=2em]
    \item \textbf{Upgrade basic verbs:} good → beneficial / enjoyable; bad → harmful / negative; help → assist / boost / contribute to.
    \item \textbf{Use formal linkers:} but → however / although; so → therefore / consequently.
    \item \textbf{Nominalisation:} study hard → consistent study; do exercise → engage in physical activity.
\end{itemize}


\section{Oral English}

\begin{table}[h]
\centering
\caption{Common English Oral Mistakes \& Corrections}
\label{tab:oral_mistakes}
\begin{tabular}{|c|p{4cm}|p{4cm}|c|p{4cm}|}
\hline
\textbf{No.} & \textbf{Wrong Expression ❌} & \textbf{Right Expression ✅} & \textbf{Error Type} & \textbf{Deep Analysis \& Memory Anchor} \\ \hline
1 & I very like this movie. & I am really into this movie. / I highly recommend it. & Collocation & \textbf{CN Trap}: "Very" cannot directly modify verbs like "like". \newline \textbf{Upgrade}: Use "be obsessed with" (addicted) or "it appeals to me" (attractive). \\ \hline
2 & Because AI is convenient, so we use it. & Because AI is convenient, we use it. / AI is convenient, so we use it. & Conjunction Logic & \textbf{Grammar Law}: "Because" and "So" cannot coexist in one sentence. \newline \textbf{Bonus}: "Due to the convenience of AI" for formal writing. \\ \hline
3 & The price of this book is expensive. & The book is expensive. / The price is high. & Subject-Verb Logic & \textbf{Physics}: Prices have "High/Low"; Items have "Expensive/Cheap". \newline \textbf{Native}: "It costs a fortune." \\ \hline
4 & With the development of society, ... & As society evolves, ... / With rapid technological advancement, ... & Template Phrasing & \textbf{Avoid}: This phrase sounds childish in PET/FCE exams. \newline \textbf{Tip}: Cut to the chase or use "As... unfolds". \\ \hline
5 & I think this is a good thing. & From my perspective, this brings more pros than cons. / It’s a double-edged sword. & Opinion Stating & \textbf{Upgrade}: Avoid simple "Good/Bad". \newline \textbf{Logic}: Use "Admittedly", "However", "Ultimately" to build a logical chain. \\ \hline
\end{tabular}
\end{table}



\end{document}
